\subsection{NORMED SPACES OVER A VALUED DIVISION RING}

DEFINITION 5. If E is a (left) vector space over a non-discrete valued division ring K, a norm on E is a mapping \(\mathbf{x} \to p(\mathbf{x})\) of E into \(\mathbf{R}_+\) which satisfies the following axioms:

\[
\left(\mathrm{NO} _ {\mathbf {I I}}\right) \quad p (\boldsymbol {x} + \boldsymbol {y}) \leqslant p (\boldsymbol {x}) + p (\boldsymbol {y}) \text {   for   all   } \boldsymbol {x}, \boldsymbol {y} \text {   in   } \mathbf {E};
\]

\[
\left(\mathrm{NO} _ {\mathbf {I I I}}\right) \quad p (t \boldsymbol {x}) = | t | p (\boldsymbol {x}) \text {for all} t \in \mathrm{K} \text {and all} \boldsymbol {x} \in \mathrm{E}.
\]

The normed spaces most frequently met with have either R or C as field of scalars (with the usual absolute value).

From \((\mathrm{NO}_{\mathrm{III}})\) it follows in particular that \(p(-x)=p(x)\) ; hence if we put \(d(x,y)=p(x-y)\) , d is an invariant metric on the additive group E, and defines a metric space topology compatible with the additive group structure of E (no. 1, Proposition 3); moreover, the mapping

\[
(t, \boldsymbol {x}) \rightarrow t \boldsymbol {x}
\]

is continuous on \(\mathbf{K} \times \mathbf{E}\) ; for we have

\[
t \boldsymbol {x} - t _ {0} \boldsymbol {x} _ {0} = (t - t _ {0}) (\boldsymbol {x} - \boldsymbol {x} _ {0}) + (t - t _ {0}) \boldsymbol {x} + t _ {0} (\boldsymbol {x} - \boldsymbol {x} _ {0})
\]

and therefore

\[
p (t \boldsymbol {x} - t _ {0} \boldsymbol {x} _ {0}) \leqslant | t - t _ {0} | p (\boldsymbol {x} - \boldsymbol {x} _ {0}) + | t - t _ {0} | p (\boldsymbol {x}) + | t _ {0} | p (\boldsymbol {x} - \boldsymbol {x} _ {0}),
\]

which shows that the left-hand side can be made as small as we please by taking \(|t - t_{0}|\) and \(p(x - x_{0})\) sufficiently small.

DEFINITION 6. If K is a non-discrete valued division ring, a vector space E over K, endowed with the structure defined by a given norm on E, is called a normed space over K.

A normed space will always be considered as endowed with the topology and the uniformity defined by its norm.

Examples. 1) On a non-discrete valued division ring K, considered as a (left or right) vector space over itself, the absolute value \(|x|\) is a norm.

\begin{enumerate}
\def\labelenumi{\arabic{enumi})}
\setcounter{enumi}{1}
\item
  The expression \(||x|| = \sqrt{\sum_{i=1}^{n} x_i^2}\) , which we have called the Euclidean norm on the space \(\mathbf{R}^n\) (Chapter VI, § 2) is evidently a norm in the sense of Definition 5. So are the functions \(\sup_{1 \leqslant i \leqslant n} |x_i|\) and \(\sum_{i=1}^{n} |x_i|\) .
\item
  Let \(\mathcal{B}(\mathrm{E};\mathrm{K})\) be the set of all functions \(f\) on a set \(\mathrm{E}\) which take their values in a non-discrete valued division ring \(\mathbf{K}\) and are such that the real-valued function \(x\to |f(x)|\) is bounded on \(\mathbf{E}\) . This set is clearly a vector subspace of the (left or right) vector space \(\mathbf{K}^{\mathbf{E}}\) of all mappings of \(\mathbf{E}\) into \(\mathbf{K}\) . If we put \(p(f)=\sup_{x\in\mathbb{E}}|f(x)|\) , then \(p\) is a norm on the vector space \(\mathcal{B}(\mathrm{E};\mathrm{K})\) (cf.~Chapter X, § 1).
\end{enumerate}

* 4) On the vector space \(\mathcal{C}(\mathrm{I};\mathbf{R})\) of all finite continuous real-valued functions defined on the interval \(\mathbf{I} = [0,1]\) of \(\mathbf{R}\) , the function

\[
p (x) = \int_ {0} ^ {1} | x (t) | d t
\]

is a norm. *

In a normed space E, the (closed) ball B with centre o and radius i, that is to say the set of all \(x \in E\) such that \(p(x) \leqslant i\) , will be called the unit ball in E. Let us show that a fundamental system of neighbourhoods of o in E is formed by the transforms of the unit ball by the homotheties \(x \to tx\) , where t runs through the set of non-zero elements of K. The image of B under this homothety is the closed ball with centre o and radius \(|t|\) ; hence it is enough to show that for each real number r \textgreater{} o there exists \(t \in K\) such that \(o < |t| < r\) . Now since the absolute value of K is not improper, there exists \(t_{0} \in K\) such that \(o < |t_{0}| < i\) ; it is therefore enough to take \(t = t_{0}^{n}\) , where n is a sufficiently large integer, in order that \(|t| = |t_{0}|^{n} < r\) .

DEFINITION 7. Two norms on a vector space E (over a non-discrete valued division ring K) are said to be equivalent if they define the same topology on E.

PROPOSITION 7. Two norms \(p, q\) on a vector space \(E\) are equivalent if and only if there exist two numbers \(a > 0, b > 0\) such that

\[
a. p (\pmb {x}) \leqslant q (\pmb {x}) \leqslant b. p (\pmb {x})\tag{I}
\]

for all \(x \in E\) .

These inequalities are sufficient, for it follows from the relation \(a.p(x) \leqslant q(x)\) that, for each \(r > 0\) , the closed ball with centre \(o\) and radius \(ar\) (relative to the norm \(q\) ) is contained in the closed ball with centre \(o\) and radius \(r\) (relative to the norm \(p\) ); hence the topology defined by \(q\) is finer than the topology defined by \(p\) . Similarly the inequality \(q(x) \leqslant b.p(x)\) shows that the topology defined by \(p\) is finer than that defined by \(q\) , and hence \(p\) and \(q\) are equivalent.

Let us now show that the inequalities (1) are necessary. If the topology defined by \(q\) is finer than the topology defined by \(p\) , then the unit ball with respect to \(p\) contains a closed ball with centre \(o\) and radius \(\alpha > 0\) with respect to \(q\) ; i.e., the relation \(q(x) \leqslant \alpha\) implies \(p(x) \leqslant 1\) . If \(t_0 \in K\) is such that \(0 < |t_0| < 1\) , then for each \(x \neq 0\) in \(E\) there is a unique rational integer \(k\) such that \(\alpha |t_0| < q(t_0^k x) \leqslant \alpha\) ; therefore \(p(t_0^k x) \leqslant 1\) , so that

\[
p (\boldsymbol {x}) \leqslant \frac {\mathrm{I}}{\left| t _ {0} \right| ^ {k}} \leqslant \frac {\mathrm{I}}{\alpha \left| t _ {0} \right|} q (\boldsymbol {x});
\]

putting \(a = \alpha |t_0|\) we have therefore \(a.p(x) \leqslant q(x)\) for all \(x \neq 0\) , and this inequality is also valid when \(x = 0\) . Similarly we show that if the topology defined by \(p\) is finer than that defined by \(q\) , there exists \(b > 0\) such that \(q(x) \leqslant b.p(x)\) .

Example. In the space \(R^{n}\) , the three norms

\[
\sqrt {\sum_ {i = 1} ^ {n} x _ {i} ^ {2}}, \quad \sup _ {1 \leqslant i \leqslant n} | x _ {i} | \quad \text { and } \quad \sum_ {i = 1} ^ {n} | x _ {i} |
\]

are equivalent, because we have

\[
\sup _ {1 \leqslant i \leqslant n} | x _ {i} | \leqslant \sqrt {\sum_ {i = 1} ^ {n} x _ {i} ^ {2}} \leqslant \sum_ {i = 1} ^ {n} | x | _ {i} \leqslant n. \sup _ {1 \leqslant i \leqslant n} | x _ {i} |.\tag{2}
\]

PROPOSITION 8. Let E be a normed space over a non-discrete valued division ring, let p be the norm on E, and let \(\hat{E}\) be the additive topological group which is the completion of the additive group E. Then the function \((t, x) \to tx\) can be extended by continuity to \(\hat{K} \times \hat{E}\) and defines on \(\hat{E}\) a vector space structure over \(\hat{K}\) ; the norm p can be extended by continuity to a norm \(\overline{p}\) on \(\hat{E}\) which defines the topology of \(\hat{E}\) .

The extension of tx by continuity is a particular case of the theorem of extension of a continuous bilinear mapping of a product of two abelian groups into a third (Chapter III, § 6, no. 5, Theorem 1); we have \(\mathbf{I} \cdot \mathbf{x} = \mathbf{x}\) and \(t(ux) = (tu)x\) for \(t \in \hat{\mathbf{K}}\) , \(u \in \hat{\mathbf{K}}\) and \(\mathbf{x} \in \hat{\mathbf{E}}\) , by the principle of extension of identities; hence the external law \((t, \mathbf{x}) \to tx\) indeed defines on \(\hat{\mathbf{E}}\) a structure of a vector space over \(\hat{\mathbf{K}}\) . On the other hand, the invariant metric \(\overline{d}(\mathbf{x}, \mathbf{y}) = \overline{p}(\mathbf{x} - \mathbf{y})\) extends to an invariant metric \(\overline{d}\) on \(\hat{\mathbf{E}}\) ( \(\S 2\) , no. 1, Proposition 1) which defines the topology of \(\hat{\mathbf{E}}\) ; if we set \(\overline{p}(\mathbf{x}) = \overline{d}(\mathrm{o}, \mathbf{x})\) , then \(\overline{p}\) is the extension of \(p\) by continuity, and satisfies axioms \((\mathrm{NO}_{\mathbf{I}})\) and \((\mathrm{NO}_{\mathbf{II}})\) ; by virtue of the continuity of \(tx\) on \(\hat{\mathbf{K}} \times \hat{\mathbf{E}}\) , \(\overline{p}\) also satisfies \((\mathrm{NO}_{\mathbf{III}})\) (principle of extension of identities) and is therefore a norm on \(\hat{\mathbf{E}}\) .

When we have to consider a definite normed space structure on a vector space E, we shall usually denote the norm of a vector x by \(\|x\|\) , unless this notation is likely to lead to confusion.

